%! TEX root = master.tex
\chapter{1}{Coordinate Proofs}

A coordinate proof of a statement in plane geometry has three steps. First,
place the figure on the axes where the algebra is simplest, without losing
generality. Then name the points with as few letters as possible. Finally,
turn lengths, midpoints, parallelism and perpendicularity into algebra. Most
of the care goes into the first step: a placement that quietly makes the
figure special proves only a special case.

\section{Distance and division of a segment}\label{sec:distance}

%! TEX root = ../master.tex
\begin{theorem}[Distance formula]\label{thm:distance-formula}
  The distance between $P_1(x_1, y_1)$ and $P_2(x_2, y_2)$ is
  \[
    P_1P_2 = \sqrt{(x_2 - x_1)^2 + (y_2 - y_1)^2}.
  \]
\end{theorem}

%! TEX root = ../master.tex
\begin{theorem}[Division of a segment]\label{thm:section-formula}
  Let $P_1(x_1, y_1)$ and $P_2(x_2, y_2)$ be distinct points.
  \begin{enumerate}[label=(\alph*)]
    \item The point $P$ on $P_1P_2$ with $P_1P : PP_2 = r : s$ (where
      $r, s > 0$) is
      \[
        P = \left( \frac{s x_1 + r x_2}{r + s}, \frac{s y_1 + r y_2}{r + s} \right).
      \]
      In particular the midpoint is
      $\bigl( \frac{x_1 + x_2}{2}, \frac{y_1 + y_2}{2} \bigr)$.
    \item If $P(x, y)$ and $Q(x', y')$ also lie on the line $P_1P_2$, then
      \[
        \frac{PQ}{P_1P_2} = \frac{\abs{x' - x}}{\abs{x_2 - x_1}}
        \quad\text{(line not vertical)}, \qquad
        \frac{PQ}{P_1P_2} = \frac{\abs{y' - y}}{\abs{y_2 - y_1}}
        \quad\text{(line not horizontal)}.
      \]
  \end{enumerate}
\end{theorem}

\begin{proof}
  Every point of the line is $P_1 + t\,(x_2 - x_1, y_2 - y_1)$ for exactly one
  real $t$, and by \cref{thm:distance-formula} the points with parameters $t$
  and $t'$ are $\abs{t' - t} \cdot P_1P_2$ apart. Their $x$-coordinates differ
  by $\abs{t' - t} \abs{x_2 - x_1}$ and their $y$-coordinates by
  $\abs{t' - t} \abs{y_2 - y_1}$, which gives (b). For (a), $P$ lies between
  $P_1$ and $P_2$ at $t = r / (r + s)$, and
  $x_1 + \frac{r}{r + s}(x_2 - x_1) = \frac{s x_1 + r x_2}{r + s}$, and likewise
  for $y$.
\end{proof}


\begin{exercise}[Varignon]\label{ex:varignon}
  Prove that joining the midpoints of the sides of a quadrilateral, in
  order, gives a parallelogram.
\end{exercise}

\begin{figure}[htbp]
  \centering
  \includegraphics[width=0.72\linewidth]{midpoint-quadrilateral}
  \caption{The quadrilateral $ABCD$ and its midpoint quadrilateral, with
    sides $j$, $k$, $l$, $m$.}
  \label{fig:midpoint-quadrilateral}
\end{figure}

\begin{proof}
  Place $A$ at the origin and $B$ on the positive $x$-axis, so that
  $A(0, 0)$, $B(a, 0)$, $C(b, c)$ and $D(d, e)$, as in
  \cref{fig:midpoint-quadrilateral}. Any quadrilateral can be moved into this
  position. By \cref{thm:section-formula}, the midpoints of $AB$, $BC$, $CD$
  and $DA$ are
  \[
    K\Bigl(\frac a2, 0\Bigr),\quad
    L\Bigl(\frac{a + b}{2}, \frac c2\Bigr),\quad
    M\Bigl(\frac{b + d}{2}, \frac{c + e}{2}\Bigr),\quad
    N\Bigl(\frac d2, \frac e2\Bigr).
  \]
  Compare opposite sides of $KLMN$ by their changes in $x$ and $y$:
  \[
    \begin{aligned}
      K \to L &: \Bigl(\frac b2, \frac c2\Bigr), &
      N \to M &: \Bigl(\frac b2, \frac c2\Bigr), \\
      K \to N &: \Bigl(\frac{d - a}{2}, \frac e2\Bigr), &
      L \to M &: \Bigl(\frac{d - a}{2}, \frac e2\Bigr).
    \end{aligned}
  \]
  Each pair of opposite sides has the same change in $x$ and the same change
  in $y$. By \cref{thm:distance-formula}, opposite sides therefore have
  equal length: $k = m = \frac12\sqrt{b^2 + c^2}$ and
  $j = l = \frac12\sqrt{(d - a)^2 + e^2}$. They are also parallel, since
  they have equal slopes or are both vertical. A quadrilateral whose
  opposite sides are parallel is a parallelogram.
\end{proof}

\begin{remark}
  The proof shows more. $KL$ and $NM$ both have the same changes as
  $\frac12 \overrightarrow{AC} = (\frac b2, \frac c2)$, so each is parallel to
  the diagonal $AC$ and half its length. Likewise $KN$ and $LM$ are parallel
  to $BD$ and half its length.
\end{remark}

\begin{exercise}\label{ex:trisection}
  Let $M$ be the midpoint of the side $AB$ of the parallelogram $ABCD$.
  Prove that the segment $MC$ and the diagonal $BD$ trisect each other.
\end{exercise}

\begin{figure}[htbp]
  \centering
  \includegraphics[width=0.48\linewidth]{parallelogram-trisection}\hfill
  \includegraphics[width=0.5\linewidth]{trisection-ratio}
  \caption{The parallelogram with $A$ at the origin and $D$ on the $x$-axis.
    Left: the unknown intersection $P$. Right: $P$ found.}
  \label{fig:trisection}
\end{figure}

\begin{proof}
  Place the parallelogram as in \cref{fig:trisection}: $A(0, 0)$,
  $B(2a, 2b)$, $D(c - 2a, 0)$ and $C(c, 2b)$, with $b > 0$ and $c \neq 2a$.
  Then $\overrightarrow{AB} = \overrightarrow{DC}$, so $ABCD$ really is a
  parallelogram. The factors of $2$ make the midpoint $M(a, b)$ free of
  fractions.

  \emph{Finding $P$.} Suppose first that $c \neq a$ and $c \neq 4a$, so
  neither line is vertical. Then
  \[
    MC : y = b + \frac{b}{c - a}(x - a), \qquad
    BD : y = 2b + \frac{2b}{4a - c}(x - 2a).
  \]
  Set them equal, divide by $b$, and rearrange:
  \[
    \frac{x - a}{c - a} - \frac{2(x - 2a)}{4a - c} = 1.
  \]
  Multiplying by $(c - a)(4a - c)$ gives
  \[
    (x - a)(4a - c) - 2(x - 2a)(c - a) = (c - a)(4a - c),
  \]
  which expands to
  \[
    3(2a - c)\,x - 8a^2 + 5ac = 5ac - c^2 - 4a^2.
  \]
  So $3(2a - c)\,x = 4a^2 - c^2 = (2a - c)(2a + c)$. Since $c \neq 2a$,
  \[
    x = \frac{2a + c}{3}.
  \]
  Then $x - a = \frac{c - a}{3}$, and the equation of $MC$ gives
  \[
    y = b + \frac{b}{c - a} \cdot \frac{c - a}{3} = \frac{4b}{3}.
  \]
  So the lines meet at $P\bigl( \frac{2a + c}{3}, \frac{4b}{3} \bigr)$.

  \emph{Ratios.} Neither segment is horizontal, so by
  \cref{thm:section-formula}(b) lengths along each one are in the same ratio
  as the changes in $y$. Along $BD$, the $y$-coordinates of $B$, $P$, $D$ are
  $2b$, $\frac{4b}{3}$, $0$, so
  \[
    BP : PD = \frac{2b}{3} : \frac{4b}{3} = 1 : 2.
  \]
  Along $MC$, the $y$-coordinates of $M$, $P$, $C$ are $b$, $\frac{4b}{3}$,
  $2b$, so
  \[
    MP : PC = \frac{b}{3} : \frac{2b}{3} = 1 : 2.
  \]
  So $P$ is a point of trisection of both segments: each segment cuts the
  other one-third of the way along.
\end{proof}

\begin{remark}
  \Cref{thm:section-formula}(a) gives a shorter proof that also covers the
  cases $c = a$ and $c = 4a$ excluded above. The point one-third of the way
  from $M$ to $C$ is
  $\frac{2M + C}{3} = \bigl( \frac{2a + c}{3}, \frac{4b}{3} \bigr)$. The
  point one-third of the way from $B$ to $D$ is
  $\frac{2B + D}{3} = \bigl( \frac{2a + c}{3}, \frac{4b}{3} \bigr)$. They are
  the same point, so the segments meet there and trisect each other. No line
  equations are needed.
\end{remark}

\begin{exercise}\label{ex:trapezoid-midline}
  Prove that the segment joining the midpoints of the non-parallel sides of a
  trapezoid is parallel to the parallel sides, and that its length is half
  the sum of their lengths.
\end{exercise}

\begin{figure}[htbp]
  \centering
  \includegraphics[width=0.6\linewidth]{trapezoid-midline}
  \caption{A trapezoid with $A(-a, 0)$, $B(a, 0)$ and the top side at
    height $2c$. The figure shows the symmetric case $D(-b, 2c)$; the proof
    places $D$ at $(d, 2c)$.}
  \label{fig:trapezoid}
\end{figure}

\begin{proof}
  Put the parallel sides on horizontal lines with $AB$ centred at the
  origin: $A(-a, 0)$, $B(a, 0)$, $C(b, 2c)$ and $D(d, 2c)$, where $a, c > 0$
  and $d < b$. The non-parallel sides are $AD$ and $BC$, and by
  \cref{thm:section-formula} their midpoints are
  \[
    M_1\Bigl(\frac{d - a}{2}, c\Bigr), \qquad
    M_2\Bigl(\frac{a + b}{2}, c\Bigr).
  \]
  \emph{Parallel.} $AB$, $DC$ and $M_1M_2$ all join two points with equal
  $y$-coordinates, so all three have slope $0$ and are parallel.

  \emph{Length.} On horizontal segments the length is the change in $x$:
  \[
    M_1M_2 = \frac{a + b}{2} - \frac{d - a}{2} = \frac{2a + (b - d)}{2},
    \qquad AB = 2a, \qquad DC = b - d.
  \]
  So $M_1M_2 = \frac12 (AB + DC)$.
\end{proof}

\begin{remark}
  The first version of this proof placed the top side symmetrically, at
  $D(-b, 2c)$ and $C(b, 2c)$ as in \cref{fig:trapezoid}. That makes the
  figure symmetric about the $y$-axis, so it covers only isosceles
  trapezoids. Centring \emph{one} parallel side costs nothing, but the
  other side must be left free. This is the placement trap described at the
  start of the chapter.
\end{remark}

\section{Perpendicular lines}\label{sec:perpendicular}

%! TEX root = ../master.tex
\begin{theorem}[Perpendicular lines]\label{thm:perpendicular-slopes}
  Two non-vertical lines with slopes $m_1$ and $m_2$ are perpendicular if and
  only if $m_1 m_2 = -1$. A vertical line is perpendicular exactly to the
  horizontal lines.
\end{theorem}


\begin{exercise}\label{ex:rhombus}
  Prove that the diagonals of a rhombus are perpendicular.
\end{exercise}

\begin{figure}[htbp]
  \centering
  \includegraphics[width=0.55\linewidth]{rhombus-diagonals}
  \caption{The rhombus $ABCD$ with its diagonals (red).}
  \label{fig:rhombus}
\end{figure}

\begin{proof}
  Place the rhombus as in \cref{fig:rhombus}: $A(0, 0)$, $B(a, 0)$,
  $C(a + b, d)$ and $D(b, d)$, with $a > 0$ and $d \neq 0$. This is a
  parallelogram, and it is a rhombus when $AB = AD$, that is,
  \[
    a^2 = b^2 + d^2.
  \]
  This gives $b^2 - a^2 = -d^2 \neq 0$, so $b \neq \pm a$ and neither
  diagonal is vertical. Their slopes are
  \[
    m_{AC} = \frac{d}{a + b}, \qquad m_{BD} = \frac{d}{b - a},
  \]
  and
  \[
    m_{AC}\, m_{BD} = \frac{d^2}{b^2 - a^2} = \frac{d^2}{-d^2} = -1.
  \]
  By \cref{thm:perpendicular-slopes}, the diagonals are perpendicular.
\end{proof}

\begin{exercise}\label{ex:equilateral-medians}
  Prove that the medians of an equilateral triangle are altitudes.
\end{exercise}

\begin{figure}[htbp]
  \centering
  \includegraphics[width=0.5\linewidth]{equilateral-medians}
  \caption{An equilateral triangle on the base $AB$, with its medians
    (dashed).}
  \label{fig:equilateral}
\end{figure}

\begin{proof}
  Let $A(0, 0)$, $B(2a, 0)$ with $a > 0$, and $C(x, y)$ with $y > 0$, as in
  \cref{fig:equilateral}.

  \emph{Locating $C$.} The triangle is equilateral, so
  $CA^2 = AB^2 = 4a^2$ and $CB^2 = AB^2 = 4a^2$. By
  \cref{thm:distance-formula},
  \[
    x^2 + y^2 = 4a^2, \qquad (x - 2a)^2 + y^2 = 4a^2.
  \]
  Subtracting the first equation from the second gives $-4ax + 4a^2 = 0$, so
  $x = a$. Then $y^2 = 3a^2$, and since $y > 0$, $C = (a, a\sqrt3)$.

  \emph{Midpoints.} By \cref{thm:section-formula},
  \[
    M_{AB}(a, 0), \qquad
    M_{AC}\Bigl(\frac a2, \frac{a\sqrt3}{2}\Bigr), \qquad
    M_{BC}\Bigl(\frac{3a}{2}, \frac{a\sqrt3}{2}\Bigr).
  \]
  \emph{Median from $C$.} $C$ and $M_{AB}$ both have $x$-coordinate $a$, so
  this median is the vertical line $x = a$. It is perpendicular to the
  horizontal side $AB$.

  \emph{Median from $B$.} The slopes are
  \[
    m_{AC} = \frac{a\sqrt3}{a} = \sqrt3,
    \qquad
    m_{BM_{AC}} = \frac{a\sqrt3/2}{a/2 - 2a} = -\frac{1}{\sqrt3},
  \]
  whose product is $-1$.

  \emph{Median from $A$.} The slopes are
  \[
    m_{BC} = \frac{a\sqrt3}{a - 2a} = -\sqrt3,
    \qquad
    m_{AM_{BC}} = \frac{a\sqrt3/2}{3a/2} = \frac{1}{\sqrt3},
  \]
  whose product is again $-1$.

  By \cref{thm:perpendicular-slopes}, each median is perpendicular to the
  side it meets, so each median is an altitude.
\end{proof}

\begin{exercise}[Converse of Pythagoras]\label{ex:converse-pythagoras}
  Prove that if the squares of two sides of a triangle add up to the square
  of the third side, then the triangle is a right triangle.
\end{exercise}

\begin{figure}[htbp]
  \centering
  \includegraphics[width=0.55\linewidth]{converse-pythagoras}
  \caption{The triangle with the vertex in question at the origin, tilted
    off the axes so that nothing about the angle at $A$ is assumed.}
  \label{fig:converse-pythagoras}
\end{figure}

\begin{proof}
  Name the vertices so that the hypothesis reads $AB^2 + AC^2 = BC^2$, and
  put the shared vertex at the origin: $A(0, 0)$, $B(x_1, y_1)$ and
  $C(x_2, y_2)$. By
  \cref{thm:distance-formula},
  \[
    (x_1^2 + y_1^2) + (x_2^2 + y_2^2)
    = (x_2 - x_1)^2 + (y_2 - y_1)^2
    = x_1^2 + x_2^2 + y_1^2 + y_2^2 - 2x_1x_2 - 2y_1y_2.
  \]
  Cancelling the common terms leaves
  \begin{equation}\label{eq:dot-zero}
    x_1x_2 + y_1y_2 = 0.
  \end{equation}
  We show that $AB \perp AC$.

  \emph{If $x_1x_2 \neq 0$}, neither $AB$ nor $AC$ is vertical. Dividing
  \eqref{eq:dot-zero} by $x_1x_2$ gives
  \[
    \frac{y_1}{x_1} \cdot \frac{y_2}{x_2} = -1,
  \]
  which is $m_{AB}\, m_{AC} = -1$. Apply \cref{thm:perpendicular-slopes}.

  \emph{If $x_1 = 0$}, then $B \neq A$ forces $y_1 \neq 0$, and
  \eqref{eq:dot-zero} becomes $y_1 y_2 = 0$, so $y_2 = 0$. Then $AB$ lies on
  the $y$-axis and $AC$ on the $x$-axis, so they are perpendicular. The case
  $x_2 = 0$ is the same with $B$ and $C$ exchanged.

  In every case the angle at $A$ is a right angle.
\end{proof}

\begin{remark}
  The slope argument alone breaks down when a side is vertical, which is
  why the case split is needed. In vector language, \eqref{eq:dot-zero} says
  $\overrightarrow{AB} \cdot \overrightarrow{AC} = 0$, which handles every
  case at once.
\end{remark}

\begin{exercise}\label{ex:chord-bisector}
  Prove that the perpendicular bisector of a chord of a circle passes
  through the centre.
\end{exercise}

\begin{figure}[htbp]
  \centering
  \includegraphics[width=0.42\linewidth]{chord-bisector}
  \caption{The chord $BC$ and its perpendicular bisector (red), drawn
    without assuming that it passes through $A$.}
  \label{fig:chord-bisector}
\end{figure}

\begin{proof}
  Centre the circle at the origin $A(0, 0)$ with radius $r$, and let the
  chord have endpoints $B(x_1, y_1) \neq C(x_2, y_2)$ on the circle, so
  \begin{equation}\label{eq:on-circle}
    x_1^2 + y_1^2 = r^2 = x_2^2 + y_2^2.
  \end{equation}
  The midpoint of $BC$ is
  $M\bigl(\frac{x_1 + x_2}{2}, \frac{y_1 + y_2}{2}\bigr)$.

  \emph{General case: $x_1 \neq x_2$ and $y_1 \neq y_2$.} Here
  $m_{BC} = \frac{y_2 - y_1}{x_2 - x_1}$ is neither zero nor undefined. By
  \cref{thm:perpendicular-slopes}, the perpendicular bisector has slope
  $-\frac{x_2 - x_1}{y_2 - y_1}$, so its equation is
  \[
    y - \frac{y_1 + y_2}{2}
    = -\frac{x_2 - x_1}{y_2 - y_1}\Bigl(x - \frac{x_1 + x_2}{2}\Bigr).
  \]
  The origin lies on this line exactly when
  \[
    -\frac{y_1 + y_2}{2} = \frac{(x_2 - x_1)(x_2 + x_1)}{2(y_2 - y_1)}.
  \]
  Multiplying by $-2(y_2 - y_1) \neq 0$ turns this into the equivalent
  condition
  \[
    y_2^2 - y_1^2 = -(x_2^2 - x_1^2),
    \qquad\text{that is,}\qquad
    x_2^2 + y_2^2 = x_1^2 + y_1^2.
  \]
  This is true by \eqref{eq:on-circle}. So the bisector passes through $A$.

  \emph{Horizontal chord: $y_1 = y_2$.} Then $x_1 \neq x_2$, and
  \eqref{eq:on-circle} gives $x_1^2 = x_2^2$, so $x_2 = -x_1$. The
  perpendicular bisector is the vertical line through $M$, namely
  $x = \frac{x_1 + x_2}{2} = 0$. This is the $y$-axis, which contains $A$.

  \emph{Vertical chord: $x_1 = x_2$.} In the same way $y_2 = -y_1$, and the
  bisector is $y = 0$, which contains $A$.
\end{proof}

\begin{remark}
  A slope-free proof is possible. The perpendicular bisector of $BC$ is
  exactly the set of points equidistant from $B$ and $C$, and the centre is
  such a point, since $AB = AC = r$.
\end{remark}
